\documentclass[11pt]{article}
\usepackage{amsmath,amssymb}
\title{The signed gauge quotient and the remaining nonlinear wavelength current}\author{}\date{}
\begin{document}\maketitle
\paragraph{The relative block and its boundary.}
For the two-link signed transverse trace,
its eigenvalues give the stated determinant, the Haar prefactor is
$(2\pi)^3/(2\pi^2)^2 B^{-3}=2/(\pi B^3)$, and differentiation gives
$-412/(343B)$ versus classical curvature $20/7$. This is a signed leading
quadratic test; it does not estimate the nonlinear growing-block remainder.
Here all interiors are those of the admitted complete relative block in
the exact minimum flow connected response and relative green improved connected bound calculations. The exterior, including all global holonomies, stays fixed or
moves along the declared boundary path. No periodic zero mode is inverted.

\paragraph{1. Exact critical-background gauge quotient.}
At the actual conditional ACTION minimum let $H$ be the canonical Hessian
on all free links before tree fixing. Let $D$ be the covariant infinitesimal
interior gauge action, with zero boundary gauge parameter. Criticality and
exact gauge invariance give $HD=0$. The relative forest meets every interior
vertex and is rooted at fixed boundary vertices. If $L$ selects its tree
rows, $LD$ is triangular after ordering the vertices, with diagonal adjoint
rotations of determinant one (up to orientation signs). Thus
$|\det LD|=1$. Let $E$ insert independent chord variations with tree rows zero,
and put $H_{\rm tr}=E^*HE$. Define
\[
 \Delta_0=D^*D,\qquad \Delta_1=H+DD^*.
\]
For the admitted minimum both operators are positive on their respective
relative spaces. The linear change $a=Ex+D\phi$ has absolute determinant
one. Since $HD=0$, its quadratic energy is $x^*H_{\rm tr}x$.
Integrating the gauge-fixing term $|D^*a|^2/2$ in $\phi$ gives
$(2\pi)^{g/2}/\det\Delta_0$, where $g$ is the dimension of the interior
gauge space. Comparing the resulting two Gaussian integrals proves
\[
 \boxed{\det H_{\rm tr}=\frac{\det\Delta_1}{(\det\Delta_0)^2}.}       \tag{1}
\]
This is a finite-dimensional identity at the actual critical background,
not a replacement of its compact Wilson integral by a Gaussian. Constant
canonical/Haar factors still multiply the raw Gaussian normalizer.
A different nonunit tree-coordinate Jacobian must be retained explicitly.

Consequently its signed quadratic fluctuation action is
\[
 F_{\rm quad}=\tfrac12\log\det\Delta_1-\log\det\Delta_0+\text{constant}.
                                                                    \tag{2}
\]
For either operator $A$ write
$\mathcal T(A)=\operatorname{Tr}(A^{-1}\ddot A-
 A^{-1}\dot A A^{-1}\dot A)$. The complete response is
$F_{\rm quad}''=\mathcal T(\Delta_1)/2-\mathcal T(\Delta_0)$.
It contains both second-derivative and covariance terms and the full
interior gauge determinant. Changes of right/left minimum frame contribute
their commutators and accelerations to these derivatives. Under a complete
site-gauge conjugation both determinants are invariant. In particular a
positive minimum-frame Fisher term alone is not the physical response.
At a flat reference color conjugation makes the first derivative of (2)
zero, but does not set its Hessian or its transport vertices to zero.

\paragraph{2. An exact original-law residual, not Gaussian promotion.}
Use the actual right fluctuation coordinates
$U_j=U_j^*(t)\exp(\alpha z_j)$ on every independent chord. A fixed exponential
fundamental domain can be used almost everywhere; its Haar density
$j(\alpha z)$ is independent of $t$. Equivalently use the admitted chart
and retain its already controlled complementary integral. With the former
convention write, by definition on the entire integration domain,
\[
 S_t(z)=V(t)+\tfrac12 z^*H_tz+R_t(z),\qquad
 d\mu_t=Z_t^{-1}j(\alpha z)e^{-z^*H_tz/2-R_t(z)}dz.                  \tag{3}
\]
Here $H_t=H_{{\rm tr},t}$ and $V$ is the exact classical minimum. All raw
canonical/Haar constants are independent of $t$ and remain in $Z$.
All derivatives below are at fixed RIGHT fluctuation $z$, so include the
actual minimum velocity and acceleration. Set $q=z^*\dot H z/2$ and
$F=-\log Z_{\rm original}-V$. Direct differentiation of the ORIGINAL
integral gives
\[
 F''=E_\mu[\tfrac12 z^*\ddot H z+\ddot R]
           -\operatorname{Var}_\mu(q+\dot R).                       \tag{4}
\]
Let $\gamma$ be the centered full Gaussian with covariance $H^{-1}$.
Subtracting its exact signed determinant response yields the identity
\begin{align}
 F''-F_{\rm quad}''={}&\tfrac12\operatorname{Tr}
   [\ddot H(E_\mu zz^*-H^{-1})]+E_\mu\ddot R\\
 &-[\operatorname{Var}_\mu q-\operatorname{Var}_\gamma q]
   -2\operatorname{Cov}_\mu(q,\dot R)-\operatorname{Var}_\mu\dot R.
                                                                    \tag{5}
\end{align}
Thus domain, Haar and nonlinear effects have not disappeared: they occur
in actual second and fourth moments and in the last three current terms.
A restriction to the convex chart must restore its boundary/annulus and
complement contributions to (5). The expectation-only Haar identity
$E F_J=0$ does not justify deleting acceleration terms inside $q$ or
$\ddot R$. The live nonflat baseline $F'(0)$ is still part of the stored
scalar/current pair; a centered second difference alone does not store it.

\paragraph{3. Exactly what would produce a logarithmic wavelength bound.}
A useful checkable kernel mechanism is the following finite elementary
lemma. Suppose a symmetric, componentwise kernel $\Pi(x,y)$ on a complete
four-dimensional domain has zero row sums, and
\[
 |\Pi(x,y)|\le {C\over B}(1+|x-y|)^{-6}.                            \tag{6}
\]
Suppose also that the declared boundary extension obeys the literal
translation estimate
$\sum_x|h(x+r)-h(x)|^2\le C|r|^2\|\nabla h\|_2^2$.
Then the row-sum identity rewrites its quadratic form as minus one half
the corresponding difference form. Summing shells of $O(r^3)$ sites gives
\[
 |\langle h,\Pi h\rangle|
 \le {C\log(2+R)\over B}\|\nabla h\|_2^2.                         \tag{7}
\]
The shell summand is $r^3r^{-6}r^2=r^{-1}$, not a volume factor.
If the complete physical transverse extension has the required Hodge
boundary identity, its gradient energy can be replaced by its curl energy;
for the harmonic extension this is $v^*K_{\rm eff}v$.

Neither (6) nor its row-sum/boundary completion is proved here for (2),
and (7) is not asserted for (5). Gauge Ward identities annihilate complete
pure-gauge directions; they do NOT by themselves establish the componentwise
zero-row-sum representation on a partial boundary chart. The actual
nonabelian transport vertices, second-derivative partners and interior
gauge determinant must first be combined to prove it. The actual-law
remainder (5) would then need its own bound by the SAME physical energy,
with a summable coupling correction. The relative green improved connected bound calculation only controls their
combined response by $CR^{14}/B$ times the canonical mass norm, plus its
explicit tail. It supplies no such relative curl bound.

\paragraph{Checkpoint.}
Equation (1) closes the finite critical-background gauge/normalization
bookkeeping for the proposed signed-trace route. Equation (5) isolates the
precise original-law current which cannot be replaced by a Wick estimate.
The outstanding main clause is a combined boundary/covariant kernel bound
for (2) and a wavelength-relative bound for (5). No strict coarse-scale
contraction, physical-slab occupation estimate, or S94 result follows yet.
\end{document}
